\documentclass{article}
\usepackage{enumitem}
\usepackage{amsmath}
\begin{document}

\title{Chapter 10: Cell Cycle and Cell Division}
\date{}
\maketitle

\begin{enumerate}[label=Q\arabic*.]

\item The correct sequence of phases in the cell cycle is:
\\(A) G1 $\rightarrow$ S $\rightarrow$ G2 $\rightarrow$ M $\rightarrow$ Cytokinesis
\\(B) S $\rightarrow$ G1 $\rightarrow$ G2 $\rightarrow$ M $\rightarrow$ Cytokinesis
\\(C) G1 $\rightarrow$ G2 $\rightarrow$ S $\rightarrow$ M $\rightarrow$ Cytokinesis
\\(D) M $\rightarrow$ G1 $\rightarrow$ S $\rightarrow$ G2 $\rightarrow$ Cytokinesis

\item DNA replication occurs during which phase of the cell cycle?
\\(A) G1 phase
\\(B) S phase
\\(C) G2 phase
\\(D) M phase

\item During which stage of mitosis do chromosomes line up at the metaphase plate?
\\(A) Prophase
\\(B) Metaphase
\\(C) Anaphase
\\(D) Telophase

\item The spindle checkpoint (metaphase checkpoint) ensures:
\\(A) DNA damage is repaired before replication
\\(B) All chromosomes are properly attached to spindle fibres before anaphase
\\(C) Cytokinesis occurs after karyokinesis
\\(D) Cell growth is adequate before entering S phase

\item Assertion: Meiosis results in genetically variable daughter cells.\\
Reason: Crossing over and independent assortment during meiosis I introduce new allele combinations.
\\(A) Both true, Reason is correct explanation
\\(B) Both true, Reason is not correct explanation
\\(C) Assertion true, Reason false
\\(D) Assertion false, Reason true

\item In meiosis, the reduction in chromosome number occurs at which division?
\\(A) Meiosis I (reductional division)
\\(B) Meiosis II (equational division)
\\(C) Both meiosis I and II equally
\\(D) During DNA replication before meiosis

\item Synapsis and crossing over occur during which sub-stage of prophase I?
\\(A) Leptotene
\\(B) Zygotene (synapsis) and Pachytene (crossing over)
\\(C) Diplotene
\\(D) Diakinesis

\item A diploid organism with 2n = 46 undergoes meiosis. How many chromosomes are present in each daughter cell after meiosis II?
\\(A) 46
\\(B) 23
\\(C) 92
\\(D) 12

\item In plant cells, cytokinesis occurs by:
\\(A) Cleavage furrow formation
\\(B) Cell plate formation from the center outward
\\(C) Constriction of cell membrane
\\(D) Both cleavage furrow and cell plate simultaneously

\item Which of the following statements about G0 phase is correct?
\\(A) G0 is an active phase of rapid cell division
\\(B) G0 is a quiescent state where cells exit the cell cycle temporarily or permanently
\\(C) G0 occurs between S and G2 phase
\\(D) All cells in G0 undergo apoptosis

\item The term ``karyokinesis'' refers to:
\\(A) Division of cytoplasm
\\(B) Division of nucleus
\\(C) Replication of DNA
\\(D) Formation of spindle fibres

\item During anaphase I of meiosis, what separates?
\\(A) Sister chromatids
\\(B) Homologous chromosomes
\\(C) Both sister chromatids and homologous chromosomes
\\(D) Centromeres only

\item Colchicine arrests cell division at metaphase because it:
\\(A) Inhibits DNA replication
\\(B) Depolymerizes tubulin and prevents spindle formation
\\(C) Blocks cytokinesis
\\(D) Inhibits checkpoint proteins

\item Match the following phases of meiosis I with their key events:\\
1. Leptotene $\rightarrow$ (a) Chromosomes begin to condense\\
2. Zygotene $\rightarrow$ (b) Synapsis of homologous chromosomes\\
3. Pachytene $\rightarrow$ (c) Crossing over\\
4. Diplotene $\rightarrow$ (d) Bivalents begin to separate, chiasmata visible
\\(A) 1-a, 2-b, 3-c, 4-d
\\(B) 1-b, 2-a, 3-d, 4-c
\\(C) 1-c, 2-d, 3-a, 4-b
\\(D) 1-d, 2-c, 3-b, 4-a

\item The significance of meiosis for sexual reproduction is:
\\(A) It produces identical daughter cells maintaining chromosome number
\\(B) It halves the chromosome number, maintains ploidy across generations, and introduces genetic variation
\\(C) It only occurs in somatic cells
\\(D) It produces cells with double the original DNA content

\end{enumerate}
\end{document}
